% -*- Document-Type: LaTeX; Master-File: manual.tex -*- 
% test pad
\documentclass{article}


\usepackage{graphics,epsfig,makeidx,latexsym,fancyhead,supertabular}
\usepackage{./lgrind-ck}
\LGnuminterval=0\LGindent=\parindent
%\usepackage{lgrind}
% (pas de num de ligne)

%\newenvironment{elanCode}{\par\medskip
%\begin{lgrind}\scriptsize}%
%{\end{lgrind}\smallbreak}  

%\newcommand{\insertCode}[2]{%
%\begin{elanCode} \input{#1/ElanTeX/#2} \end{elanCode}}
 
%\newcommand{\insertExamples}[1]{\insertCode{ElanExamples}{#1}}
%\newcommand{\insertLibCommon}[1]{\insertCode{ElanLibCommon}{#1}}
%\newcommand{\insertLibNoquote}[1]{\insertCode{ElanLibNotquote}{#1}}
%\newcommand{\insertLibStrategy}[1]{\insertCode{ElanLibStrategy}{#1}}

%%%%%%%%%%%%%%%%%%%%%%%%%%%%%%
\newcommand{\versionElan}{V2.00}

\newcommand{\elan}{{\sf ELAN}}
\newcommand{\EX}{\begin{example}\begin{it}}
\newcommand{\FEX}{\end{it}\end{example}}
\input macros

\begin{document}

\tableofcontents
\newpage

%%%%%%%%%%%%%%%%%%%%%%%%%%%%%%%%%%%%%%%%%%%%%%%%%%%%%%%%%%%%
\section{Syntax}
%%%%%%%%%%%%%%%%%%%%%%%%%%%%%%%%%%%%%%%%%%%%%%%%%%%%%%%%%%%%

\subsection{Lexical unities}
%%%%%%%%%%%%%%%%%%%%%%%%%%%%%%%%%%%%%%%%%%%%%%%%%%%%%%%%%%%%
\index{letter}
\grrule{letter}{\tt
A\sor B\sor C\sor D\sor E\sor F\sor G\sor H\sor 
I\sor J\sor K\sor L\sor M\sor N\sor O\sor P\sor Q\sor R\sor S\sor T\sor 
U\sor V\sor W\sor X\sor Y\sor Z 
\gror
 a\sor b\sor c\sor d\sor e\sor f\sor g\sor h\sor 
i\sor j\sor k\sor l\sor m\sor n\sor o\sor p\sor q\sor r\sor s\sor t\sor 
u\sor v\sor w\sor x\sor y\sor z
}

\index{numeral}
\grrule{numeral}{\tt
0\sor 1\sor 2\sor 3\sor 4\sor 5\sor 6\sor 7\sor 8\sor 9
}

\index{identifier}
\grrule{identifier}{ 
  \nt{letter} \itzero{\nt{letter} \sor \nt{numeral} \sor ~\sym{\_}}
}

\grrule{number}{ \itone{\nt{numeral}}}

\subsection{Definition of signatures}
%%%%%%%%%%%%%%%%%%%%%%%%%%%%%%%%%%%%%%%%%%%%%%%%%%%%%%%%%%%%

\index{sort name}\index{name!sort}
\grrule{sort name} {
  \nt{identifier}
 \gror
  \nt{identifier}\sym{[}\nt{sort name}\itzero{\sym{,}\nt{sort name}}\sym{]}
}

\index{symbol declaration}\index{declaration!function}
\grrule{symbol declaration}{
  \nt{new symbol declaration}
\gror
  \nt{alias}
}

\grrule{new symbol declaration}{
   \nt{identifier} \sym{:} \nt{rank} \omm{\nt{options}}
}

\index{alias}
\grrule{alias}{\nt{new symbol declaration} \rw{alias} \nt{old name}\sym{:}}


\grrule{name} {\itone{\nt{symbol}}}
\grrule{symbol}{
  \nt{single lexem}
\gror
  \sym{@}
\gror
  \sym{'}\nt{single lexem}\sym{'}
\gror
  \sym{'$\backslash$}\nt{number}\sym{'}
}
\grrule{single lexem}{
  \nt{identifier}
\gror
  \nt{number}
\gror
  \nt{nchar}
}

\index{selector}
\grrule{rank}{
  \nt{sort name}
\gror
  \sym{(} \itone{\nt{sort name}} \sym{)} \nt{sort name}
\gror
  \sym{(} \itone{\nt{name} \sym{:} \nt{sort name}} \sym{)} \nt{sort name}
}

\index{assocLeft}\index{assocRight}\index{pri}\index{AC}\index{code}
\grrule{option}{
  \rw{assocLeft}
\gror
  \rw{assocRight}
\gror
  \rw{pri} \nt{number}
\gror
  \sym{(} \rw{AC} \sym{)}
\gror
  \rw{code} \nt{number}
}


\subsection{Definition of module}
%%%%%%%%%%%%%%%%%%%%%%%%%%%%%%%%%%%%%%%%%%%%%%%%%%%%%%%%%%%%

\grrule{module}{
        \rw{module } \nt{formal module name} \nl
        \omm{\nt{imports}}              \nl
        \omm{\nt{sort definition}}      \nl
        \omm{\nt{operator definition}}  \nl
        \itzero{\nt{family of rules}}   \nl
        \itzero{\nt{family of strategies}}      \nl
        \rw{end}
}

\grrule{imports}{
   \rw{import}                          \nlp{5}
        \omm{\rw{global} \itone{\nt{module name}} \sym{;}} \nlp{5}
        \omm{\rw{local} \itone{\nt{module name}} \sym{;}} \nl
   \rw{end}
}


\index{sort name}\index{name!sort}
\grrule{sort name} {
  \nt{identifier}
 \gror
  \nt{identifier}\sym{[}\nt{identifier}\itzero{\sym{,}\nt{identifier}}\sym{]}
}

\index{sort definition}\index{name!sort}
\grrule{sort definition} {
  \rw{sort}
  \itone{\nt{sort name}} \sym{;}
  \rw{end}
}
\index{operator definition}\index{operator!signature}
\grrule{operator definition}{
          \rw{operators}                \nlp{5}
          \omm{\rw{global} \itone{\nt{symbol declaration}\sym{;}}}      \nlp{5}
          \omm{\rw{local} \itone{\nt{symbol declaration}\sym{;}}}       \nl
          \rw{end}
}

\subsection{Definition of rules}
%%%%%%%%%%%%%%%%%%%%%%%%%%%%%%%%%%%%%%%%%%%%%%%%%%%%%%%%%%%%

\grrule{family of rules}{
        \rw{rules for} \nt{sort name}                \nl
        \omm{\itone{\nt{variable declare}}}           \nl
        \omm{\rw{global} \itone{\nt{labelled or unlabelled rule}}} \nl
        \omm{\rw{local}  \itone{\nt{labelled rule}}} \nl
        \rw{end}
}

\grrule{labelled or unlabelled rule}{
        \nt{labelled rule}\nl
\gror   \nt{labelled rule}
}

\grrule{labelled rule}{
        \sym{[} \nt{rule name} \sym{]} \nt{rule body} \rw{end}
}

\grrule{unlabelled rule}{
        \sym{[} \sym{]} \nt{rule body} \rw{end}
}

\grrule{rule body}{
        \nt{term} \sym{=$>$} \nt{rhs}\
}

\grrule{rhs}{
        \nt{term}
\gror   \nt{rhs} \rw{if} \nt{boolean term}    \nl
\gror   \nt{rhs} \rw{where} \nt{var name} \sym{:=}
                 \sym{(} \omm{\nt{strategy}} \sym{)} \nt{term} \nl
\gror   \nt{rhs} \rw{where} \nt{var name} \sym{:=}
                 \sym{[} \nt{Pstrategy} \sym{]} \nt{term} \nl
\gror   \nt{rhs} \rw{where} \sym{(} \nt{sort} \sym{)} \nt{term} \sym{:=}
                 \sym{(} \omm{\nt{strategy}} \sym{)} \nt{term} \nl
\gror   \nt{rhs} \rw{where} \sym{(} \nt{sort} \sym{)} \nt{term} \sym{:=}
                 \sym{[} \nt{Pstrategy} \sym{]} \nt{term} \nlp{2}
\gror   \rw{switch}\nlp{5}
        \itone{\rw{case} \nt{boolean term} \rw{then} \nt{rhs} }\nlp{5}
        \omm{\rw{otherwise} \nt{rhs}} \nl
        \rw{end}
}


\subsection{Definition of strategies}
%%%%%%%%%%%%%%%%%%%%%%%%%%%%%%%%%%%%%%%%%%%%%%%%%%%%%%%%%%%%

\grrule{family of strategies}{
        \rw{strategies for} \nt{sort name}                \nl
        \omm{\rw{global} \itone{\nt{non complete strategy}}} \nl
        \omm{\rw{local}  \itone{\nt{non complete strategy}}} \nl
        \rw{end}
}

\grrule{non complete strategy}{
        \sym{[} \nt{strategy name} \sym{]} \nt{strategy} \rw{end}\nl
}

\grrule{strategy}{
       \nt{strategy name}
\gror  \nt{choosing}
\gror  \nt{concatenation}
\gror  \nt{iterator}
\gror  \nt{normalise}
\gror  \rw{id}
}

\grrule{strategy name}{\nt{identifier}}

\grrule{strategy\_rule}{
       \nt{strategy}
\gror  \nt{rule name}
}

\grrule{choosing}{
       \rw{dont care choose} \sym{(} \nt{strategy\_rule} 
                \itzero{\sym{$||$} \nt{strategy\_rule}} \sym{)}
\gror  \rw{dont know choose} \sym{(} \nt{strategy\_rule} 
                \itzero{\sym{$||$} \nt{strategy\_rule}} \sym{)}
\gror  \rw{dc} \sym{(} \nt{strategy\_rule} 
                \itzero{\sym{$||$} \nt{strategy\_rule}} \sym{)}
\gror  \rw{dk} \sym{(} \nt{strategy\_rule} 
                \itzero{\sym{$||$} \nt{strategy\_rule}} \sym{)}
}

\grrule{concatenation}
  {\nt{strategy} \sym{;} \nt{strategy}}

\grrule{iterator}{
       \rw{iterate} \sym{(} \nt{strategy} \sym{)} 
\gror  \rw{repeat*} \sym{(} \nt{strategy} \sym{)} 
\gror  \rw{repeat+} \sym{(} \nt{strategy} \sym{)} 
}

\grrule{normalise}{
       \rw{normalise} \sym{(} \nt{strategy} \sym{)} 
}

\subsection{How to transform a program ?}
%%%%%%%%%%%%%%%%%%%%%%%%%%%%%%%%%%%%%%%%%%%%%%%%%%%%%%%%%%%%
You have to modify declaration parts:
\begin{itemize}
\item \texttt{import ...} becomes \texttt{import ... end}
\item \texttt{sort ...} becomes \texttt{sort ... end}
\item \texttt{op ... endop} becomes \texttt{operators ... end}
\end{itemize}

Now, there is only one constructions to define rules, you have to
replace: 
\begin{itemize}
\item \texttt{rule for ...} by \texttt{rules for ...}
\item \texttt{rule <name> for ...} by \texttt{rules for ... [name] ...}
\item suppress \texttt{declare} keyword
\item \texttt{body} or \texttt{bodies} by \texttt{global} or
  \texttt{local}
\item texttt{end of rule}, texttt{end of rules} by \texttt{end}
\end{itemize}

Definition of strategies looks like rules definitions, you have to
replace: 
\begin{itemize}
\item \texttt{strategy ...} by \texttt{strategies for ... global|local
    [name] ...} 
\item \texttt{iterate ... enditerate} by \texttt{iterate(...)}
\item \texttt{while ... endwhile} by \texttt{repeat*(...)}
\item \texttt{repeat ... endrepeat} by \texttt{repeat*(...)}
\item \texttt{dont know choose(r1 r2 r3)} by \texttt{dc(r1 || r2 ||
    r3)}
\item \texttt{repeat ... endrepeat dont care choose(...)} by
  \texttt{repeat*(...) ; dc(...)}
\item \texttt{end of strategy} by \texttt{end}
\item \texttt{end of module} by \texttt{end}
\end{itemize}

Local affectations and conditions are now evaluated from top to
bottom. You have to reverse the order:

\begin{verbatim}
  [] l => r
          if c2
          where v2:=(...) ...
          where v1:=(...) ...
          if c1
\end{verbatim}

becomes in the new syntax:
\begin{verbatim}
  [] l => r
          if c1
          where v1:=(...) ...
          where v2:=(...) ...
          if c2
\end{verbatim}

%%%%%%%%%%%%%%%%%%%%%%%%%%%%%%%%%%%%%%%%%%%%%%%%%%%%%%%%%%%%
\section{Library}
%%%%%%%%%%%%%%%%%%%%%%%%%%%%%%%%%%%%%%%%%%%%%%%%%%%%%%%%%%%%
The new library is split in 3 directories:
\begin{itemize}
\item \texttt{elanlib/commons} contains common modules definitions:
  int, bool, string, pair, list, array,\ldots
\item \texttt{elanlib/noquote} contains modules connected to
  identifier representation: identifier, term, substitution,
  constraint, unification,\ldots
\item \texttt{elanlib/strategy} contains modules for user-defined
  strategies: meta, any, strapply, \ldots
\end{itemize}

Definition of terms have changed, you no longer need to play with
quotes `{''}`. All modules \texttt{term}, \texttt{termF},
\texttt{termV}, \texttt{termFExt}, \texttt{termVExt} have been
replaced by only one: \texttt{termCommons} parametrized by
\texttt{Fss} and \texttt{Vars}.

\texttt{eqSystem} and \texttt{applySubstOn} are now paretrized by
\texttt{Fss} and \texttt{Vars}.

\texttt{syntacticUnification[Vars,Fss]} is replaced by
\texttt{syntacticUnification[Fss,Vars]}.
 

%%%%%%%%%%%%%%%%%%%%%%%%%%%%%%%%%%%%%%%%%%%%%%%%%%%%%%%%%%%%
\section{New Language Features}
%%%%%%%%%%%%%%%%%%%%%%%%%%%%%%%%%%%%%%%%%%%%%%%%%%%%%%%%%%%%

%----------------------------------------------------------------
% Peter's macro defs file
%\def \FULL{0} \def \ETCS{1}
\def \FULL{1} \def \ETCS{0}
%
\newcommand{\comment}[1]   {  }
\newcommand{\ELAN}{{\sf ELAN}}
%
\newcommand{\mathdef}[1]{\relax\ifmmode #1\else $#1$\fi}
\newcommand{\rrule}[2]{\mathdef{\frac{\large #1}{\large #2}}}
\renewcommand{\o}{\:\circ\:}
\newcommand{\oo}{\infty}
%
\newcommand{\te}[0]{\mathdef{\:\vdash\:}}
\newcommand{\rewrite}[0]{\mathdef{\Rightarrow}}
\newcommand{\red}[0]{_{\downarrow}}
%
\newcommand{\TT}[0]{{\cal T}}
\newcommand{\FF}[0]{{\cal F}}
\newcommand{\DD}[0]{{\cal D}}
\newcommand{\OO}[0]{{\cal O}}
\renewcommand{\SS}[0]{{\cal S}}
\newcommand{\EE}[0]{{\cal E}}
\newcommand{\XX}[0]{{\cal X}}
\newcommand{\LL}[0]{{\cal L}}
%
%\newcommand{\SS}[0]{{\cal S}}
\newcommand{\II}[0]{{\cal I}}
\newcommand{\RR}[0]{{\cal R}}
\newcommand{\CC}[0]{{\cal C}}
%
\newcommand{\IF}[0]{{\rm ~if~}}
\newcommand{\ELSE}[0]{{\rm ~else~}}
\newcommand{\FI}[0]{{\rm ~fi~}}
\newcommand{\THEN}[0]{{\rm ~then~}}
\newcommand{\WHERE}[0]{{\rm ~where~}}
\newcommand{\AND}[0]{{\rm ~and~}}
\newcommand{\OR}[0]{{\rm ~or~}}
\newcommand{\FOR}[0]{{\rm ~for~}}
\newcommand{\ANY}[0]{{\rm ~any~}}
\newcommand{\IFF}[0]{{\rm ~iff~}}
\newcommand{\IN}[0]{{\rm ~in~}}
\newcommand{\LET}[0]{{\rm ~let~}}
\newcommand{\MATCH}[0]{{< \hspace{-1mm} < ^{?}}}
%\newcommand{\Strat}[0]{{\bf Strat}}
%
\newcommand{\Strat}[1]{\langle #1 \rangle}
\newcommand{\Symb}[1]{Symbol^{#1}}
\newcommand{\Bool}[0]{Bool}
\newcommand{\List}[0]{list}
\newcommand{\appl}[2]{#1 \lceil #2 \rceil}
\newcommand{\appll}[2]{#1 \lceil \lceil #2 \rceil \rceil}
\newcommand{\MV}[0]{{\cal MV}}
%
\newcommand{\FSYM}[0]{{\sf FSYM}}
\newcommand{\RULE}[0]{{\sf LAB}}
\newcommand{\DSTR}[0]{{\sf DSTR}}
\newcommand{\DK}[0]{{\sf DK}}
\newcommand{\DC}[0]{{\sf DC}}
\newcommand{\CASE}[0]{{\sf CASE}}
\newcommand{\CONC}[0]{{\sf CONC}}
\newcommand{\ID}[0]{{\sf ID}}
\newcommand{\FIN}[0]{{\sf FIN}}
\newcommand{\FAIL}[0]{{\sf FAIL}}
%
\newcommand{\FILTER}[0]{{\sf FILTER}}
\newcommand{\dc}[0]{{\bf dc}}
\newcommand{\dk}[0]{{\bf dk}}
\newcommand{\cas}[0]{{\bf case}}
\newcommand{\nil}[0]{{\bf nil}}
\newcommand{\lab}[0]{{\bf lab}}
\newcommand{\cons}[0]{{\bf cons}}
\newcommand{\id}[0]{{\bf id}}
\newcommand{\fail}[0]{{\bf fail}}
\newcommand{\f}[0]{{\bf f}}
\newcommand{\h}[0]{{\bf h}}
%\renewcommand{\l}[0]{{\bf l}}
%\renewcommand{\r}[0]{{\bf r}}
\newcommand{\lll}[0]{{\bf l}}
\newcommand{\dd}[0]{{\bf d}}
\newcommand{\map}[0]{{\bf map}}
%
\newcommand{\while}[0]{{\bf while}}
\newcommand{\repeate}[0]{{\bf repeat}}
\newcommand{\dcwhile}[0]{{\bf dcwhile}}
\newcommand{\apply}[0]{{\bf apply}}
%
\newcommand{\joint}[0]{join}
%\newcommand{\Set}[0]{{\bf Set}}

\newcommand{\Set}[0]{{\bf List}}

\newcommand{\ter}[1]    {{\tt #1}}

\newlength{\largeurtexte}
\setlength{\largeurtexte}{134mm} % \textwidth - 5mm

\newcommand{\lban}[0]{ (  {\kern -2pt} [ }
\newcommand{\rban}[0]{ ]  {\kern -2pt} ) }
%%%%%%%%%%%%%%%%%%%%%%%%%%%%%%%%%%%%%%%%%%%%%%%%%%%%%%%%%%%%%%%%%%%%%%%%%

%----------------------------------------------------------------

\subsection{Inputs and Outputs}

Up to now, 
there exists only the following proposition, 
which is not updated for the version elan.test4.

\marginpar{
See the following library modules: prompt, io and stdio}

{\scriptsize
\subsubsection{How to read \& write}

For the beginning, let us suppose that there are two built-in symbols 
for low-level IOs (like read, write), which are defined in some
system library module (stdio.eln, io.eln, io\_aux.eln -- cf. appendix)
and handled by a built-in strategy $IO$. This assumption allows us to define the
following ELAN module for terminal IOs.
%..print
{\scriptsize
\begin{verbatim}
module print[X]
import local io[X];
op global 
  read     : X;
  print(@) : (X) X;
endop
rules for X
declare t,tt : X;
bodies
  [] read => tt     where tt:=()read(stdin)     end
  [] print(t) => tt where tt:=()write(stdout,t) end
end of rules
\end{verbatim}
}
%
The module $print$ should/could be imported with any ELAN sort $X$,
values of which we want to read in or write out. 
This module uses a system library module
$io$, which (finally) defines (the syntax of) low-level primitives 
$read$ and $write$.
By importing this module, the user can easily write out any term of any
sort. However, they can also construct IO-terms whose meanings are
not easy to find out, e.g. $print(read.read.nil)$ - supposing that (s)he
imported $print[int]$ and also $print[list[int]]$\footnote{
a behaviour depends on the ELAN reduction strategy, which is, in fact, 
left-most-inner-most}.

%...................................
\subsubsection{File Input \& Outputs}

Intuitively, the same kind of operations, as we have shown for terminal IOs,
we can do with files. The module $stdio.eln$ defines basic built-ins:
$open$ for opening and $close$ for closing files. The idea of
processes and files is that they are uniquely represented in ELAN 
by a value of the sort $Pid$. Thus, any io ($open$, $read$, $write$, $close$)
with a file identified by $pid$ goes via $pid$ (and the same for processes).

The following example shows a way how to write out a triangle into a file:
{\scriptsize
\begin{verbatim}
 10 end 
 10+ 9+ 8+ 7+ 6+ 5+ 4+ 3+ 2+ 1+ 0end 
 9+ 8+ 7+ 6+ 5+ 4+ 3+ 2+ 1+ 0end 
 8+ 7+ 6+ 5+ 4+ 3+ 2+ 1+ 0end 
 7+ 6+ 5+ 4+ 3+ 2+ 1+ 0end 
 6+ 5+ 4+ 3+ 2+ 1+ 0end 
 5+ 4+ 3+ 2+ 1+ 0end 
 4+ 3+ 2+ 1+ 0end 
 3+ 2+ 1+ 0end 
 2+ 1+ 0end 
 1+ 0end 
\end{verbatim}
}
%
The solution uses two loops ($loop$ and $loop2$) and a top-level
procedure ($wloop$), which opens and closes a file.
$loop$ writes out one line of the output file $pid$ (except the symbol $end$) and 
returns the sum of written numbers. $loop2$ writes out the whole
triangle by calling $loop$ $n$-times.
%
{\scriptsize
\begin{verbatim}
import global string io[string] io[int];  // to be able to write/read strings and ints
op global
  wloop(@,@)    : (string int) int;
  loop(@,@)     : (Pid int) int;
  loop2(@,@)    : (Pid int) int;
endop
rules for int
declare s1, s2, n, n1, n2, i    : int;
        pid, pid1               : Pid;
        file, id1               : string;
bodies
  [] loop2(pid,0) => 0                                end
  [] loop2(pid,n) => s1+s2
                 where s2 :=()loop2(pid,n-1)
                 where id1:=()write(pid,"0 end\n")             
                 where s1 :=()loop(pid,n)             end
  [] loop(pid,0) => 0                                 end
  [] loop(pid,n) => n+loop(pid,n-1) 
                 where id:=()write(pid,"+")
                 where n1:=()write(pid,n)           end
  [] wloop(file,n) => n1 if iserr(pid1) == false      // test the correctness
                 where pid1:=()close(pid)           // close the file
                 where n1  :=()loop2(pid,n)           // write out whole triangle
                 where id1 :=()write(pid,"end\n")
                 where n2  :=()write(pid,n)         // write out the first line
                 if iserr(pid) == false               // test if is is open correctly
                 where pid :=()open(file,"w")       // open file for write
  end
\end{verbatim}
}
%
As usually, $~where$'s and $if$'s in ELAN should be read from bottom to
up. In this sens, $wloop("foo",10)$ opens the file "foo" identified by
the string variable $file$,
tests whether it was correctly opened, writes out the first line 
with $10~end$, writes out the whole triangle by calling $loop2$, 
and closes the file. Any error of the $open$ or $close$ operation is detected
by $pid$ of this file.

\bigskip

The following example tries to read the file produced by $wloop$.
{\scriptsize
\begin{verbatim}
op global
  rloop(@)      : (string) int;
  sum(@,@)      : (int Pid) int;
endop
rules for int
declare s1, s2, n, n1, n2, i    : int;
        pid, pid1               : Pid;
        file, id1               : string;
bodies
  [] rloop(file) => n1  if iserr(pid1) == false             // and test
                        where pid1:=()close(pid)          // close
                        where n2  :=()write(stdout,n1)    // write out the total sum of all numbers
                        where n1  :=()sum(i,pid)            // read the triangle
                        where i   :=()read(pid)           // read a size of the triangle - the first line
                        if iserr(pid) == false              // test of correctness
                        where pid :=()open(file,"r")      // open a file for read
  end
  [] sum(0,pid) => 0                                            end
  [] sum(i,pid) => n+n2 where n2  :=()sum(i-1,pid)          // loop again
                        where n1  :=()write(stdout,n)     // write out read integer term
                        where n   :=()read(pid) if i>0    // read one line of input file    
  end
end of rules
\end{verbatim}
}
%
The body of $rloop$ called by $rloop("foo")$ opens the file "foo", reads the first line until the
symbol $end$, and it starts a loop of $sum$. What might be interesting is
that the loop of $sum$ runs $10$ times and each pass of $sum$
reads one terms of the sort $int$ until $end$, i.e. one line of
the input file. When you try to run this example for larger $n$ (like
$100$), you see that reading, analysing and normalising of 
the term $100 + 99 + \dots + 1 + 0$ takes some time.

Another (silly) example of low-level IOs tries to sketch the `interaction'
with the user. $ccode$ reads an integer (as a confidential code),
and after a `verification' ($n < 10000$) it appends this code as a log-information
into the log-file and finishes this communication.
%
{\scriptsize
\begin{verbatim}
rules for int
declare id, id1, id2            : string;
        fid, fid1               : Pid;
        n, n1, n2, m            : int;
bodies
  [] ccode => n1
        where fid1:=()close(fid)
        where id1:= ()write(stdout,"\nSorry, there are no money\n") 
        where id2:= ()write(fid,"is logged\n")    // log-info
        where n1 := ()write(fid,n)                // write code into the log-file
        where fid:= ()open(logfile,"a")           // open to append
        if n < 10000                                // validity test :-)
        where n  := ()read(stdin)
        where id := ()write(stdout,"\nEnter your code confidential: ") end
  [] ccode => 9999
        where id := ()write(stdout,"\nSorry, your code is invalid\n")   end
\end{verbatim}
}
}

%----------------------------------------------------------------
\subsection{Communication with external processes}

Up to now, 
there exists only the following proposition, 
which is not updated for the version elan.test4.

\marginpar{
See the following library modules: io and stdio}

For a reasonable example, see Castro's distributed csp-solver.

{\scriptsize
In this section we show an example of the communication between
processes (cf. Figure \ref{solvers}), namely between a solver which for some partial tasks
calls two sub-solvers ($solver1$ and $solver2$). To say a little-bit more
about these solvers (although it is not the crucial point), the solver
$solver1$ transforms an input constraint satisfaction problem $csp$
to a resolved form using rules of arch/node consistency.
For any input problem (or a term of the sort $csp$) it gives one
resolved solution (or a term of the sort $csp$).
The solver $solver2$ enumerates all solutions of a solved form, e.g.
for one solved $csp$ it may give more enumerated solutions.
The solvers $solver1$ and $solver2$ are two elan-tasks\footnote{
the only restriction to sub-solvers is that they should communicate
via streams stdin and stdout} 
started by the following script files:
%
\begin{figure}[h]
\centering
%\input{solvers.pstex_t}

\begin{picture}(0,0)%
\special{psfile=/u/protheo/borovan/elan/IOS/PROPOSITION/solvers.pstex}%
\end{picture}%
\setlength{\unitlength}{0.012500in}%
\begin{picture}(521,227)(15,570)
\end{picture}

\caption{Communication of solvers}
\label{solvers}
\end{figure}
%
{\scriptsize
\begin{verbatim}
solver1:   elan -b file1 

solver2:   elan -b file2 
\end{verbatim}
}
%
To better visualise flows of terms between processes, we can
monitor the input to $solver1$ in the Unix-file $pipe1$,
and the output in the file $pipe2$ (analogically, $pipe3$, $pipe4$ for
$solver2$):
{\scriptsize
\begin{verbatim}
tee pipe1 | elan -b file1 | tee pipe2
\end{verbatim}
}
%
Now, we can describe small piece of ELAN code which describes the 
communication of $solver$ with their two sub-solvers. It is described
in the body of $sol$:
{\scriptsize
\begin{verbatim}
module solver
import global string bool signature_csp list[csp]
              io_aux[csp,signature_csp] io[string];
op global
  sol(@)                : (csp) list[csp];
  delay(@)              : (int) int;
  read_all(@,@,@)       : (int Pid list[csp]) list[csp];
endop
rules for list[csp]
declare 
        id, id1, id2, id3, id4, id5, id6, id7, id8      : string;
        PIDA, PIDA1,PIDB,PIDB1                          : Pid;
        pb, pb1, sol, sol1, sol2                        : csp;
        n, m                                            : int;
        lst                                             : list[csp];
bodies
  [] sol(pb) => lst
          //--- kill solvers
          where PIDB1 := ()close(PIDB)
          where PIDA1 := ()close(PIDA)

          //---  read up to 100-solution of the solver2 and collect them into a list
          where lst :=()read_all(100,PIDB,nil)  
          
          //--- send the result of the solver1 to the solver2
          where id7  := ()write(stdout,"sent to solver 2\n")
          if iserror(sol1) == false and iserror(id6) == false 
          where id6   := ()write(PIDB,"end\n")
          where sol1  := ()write(PIDB,sol)

          //--- read the solution produced by the solver1
          where id2  := ()write(stdout,"received from solver1\n")
          if iserror(sol) == false
          where sol  := ()read(PIDA)

          //--- send the initial $csp$ problem to solver1
          where id1  := ()write(stdout,"sent to solver 1\n")
          if iserror(pb1) == false and iserror(id) == false 
          where id   := ()write(PIDA,"end\n")
          where pb1  := ()write(PIDA,pb)

          //--- start solver2 and test errors occured
          where id5  := ()write(stdout,"solver2 is running\n")
          if iserror(PIDB) == false 
          where PIDB  := ()create(solver2)

          //--- start solver1 and test errors occured
          where id3  := ()write(stdout,"solver1 is running\n")
          if iserror(PIDA) == false 
          where PIDA  := ()create(solver1)              
  end
end of rules
rules for list[csp]
declare 
        id, id1, id2, id3, id4, id5, id6, id7, id8      : string;
        pid                                             : Pid;
        sol, sol1                                       : csp;
        n                                               : int;
        lst                                             : list[csp];
bodies // reads all solutions up to n
  [] read_all(0,pid,lst) => lst                         end
  [] read_all(n,pid,lst) => read_all(n-1,pid,sol1.lst)
           where id1  := ()write(stdout,"\n\n")
           where sol1 := ()write(stdout,sol)
           where id   := ()write(stdout,"received solution from solver2:")
           if iserror(sol) == false
           where sol := ()read(pid)              end
  [] read_all(n,pid,lst) => lst
           where id := ()write(stdout,"\n--- FINITO ---\n\n") end
end of rules
\end{verbatim}
}
%
The following demo shows the execution of the solver $solver$. It is
necessary mentioned that between messages `sent to solver 1' and
`received from solver1' there is long delay  not because of
complexity of this problem/example, but because of loading to ELAN
with sub-tasks. If we want to execute solver for more input problems,
we should little-bit change the communication of $solver$ that
sub-solvers are loaded (created) and killed only once.
{\scriptsize
\begin{verbatim}
elan solver
 . . . . . .
enter query term finished by the key word 'end':
X1 in_domain? 1.2.3.4.5.nil.X2 in_domain?1.2.3.4.5.nil.nil,nil,
X1(+)2<=?X2.X2>=?X1(+)2.nil,nil,nil,false end

 solver1 is running 
 solver2 is running 
 sent to solver 1
 received from solver1 
 sent to solver 2
 received solution from solver2: nil,nil,nil,X1=?3.X2=?5.nil,X2>=?X1(+)2.X1(+)2<=?X2.nil,true 
 received solution from solver2: nil,nil,nil,X1=?2.X2=?5.nil,X2>=?X1(+)2.X1(+)2<=?X2.nil,true 
 received solution from solver2: nil,nil,nil,X1=?1.X2=?5.nil,X2>=?X1(+)2.X1(+)2<=?X2.nil,true 
 received solution from solver2: nil,nil,nil,X1=?3.X2=?4.nil,X2>=?X1(+)2.X1(+)2<=?X2.nil,true 
 received solution from solver2: nil,nil,nil,X1=?2.X2=?4.nil,X2>=?X1(+)2.X1(+)2<=?X2.nil,true 
 received solution from solver2: nil,nil,nil,X1=?1.X2=?4.nil,X2>=?X1(+)2.X1(+)2<=?X2.nil,true 
 received solution from solver2: nil,nil,nil,X1=?3.X2=?3.nil,X2>=?X1(+)2.X1(+)2<=?X2.nil,true 
 received solution from solver2: nil,nil,nil,X1=?2.X2=?3.nil,X2>=?X1(+)2.X1(+)2<=?X2.nil,true 
 received solution from solver2: nil,nil,nil,X1=?1.X2=?3.nil,X2>=?X1(+)2.X1(+)2<=?X2.nil,true 
 received solution from solver2: end 
\end{verbatim}
}

\comment{
\subsubsection{Library modules}

In this section, we briefly explain tree system modules which define IO-built-ins.

\subsubsubsection{stdio}

As we have said, $pid:Pid$ uniquely identifies a file or a
process. For a file, it represents a file handle, for a process it is
a process identification number. Files can be opened for write, read
or append ("w", "r", "a") which is specified by the second string of
$open$. The first one represents the name of a file. The sort $Pid$
has two constructors, which are hidden for the user, thus the user
cannot play with pids. The second $Pid$ constructor refers to a
situation where some IO-operation finished unsuccessfully. In this
case, we could (by $iserr$, and $errno$) detect an error and see its code.
%.. stdio
{\scriptsize
\begin{verbatim}
module stdio
// !!! not implemented in the compiled version !!!
import global int bool string; 
sort Pid;
op 
 global
  create(@)     : (string) Pid          code 115; // a process
  open(@,@)     : (string string) Pid   code 116; // open a file for read/write/append
  close(@)      : (Pid) Pid             code 118; // close a file/kill a process
  //------------------------------------
  stdin         : Pid;
  stdout        : Pid;
  stderr        : Pid;
  iserr(@)      : (Pid) bool;
  errno(@)      : (Pid) int;
 local
  @             : (int) Pid             code 121;
  Error(@)      : (int) Pid             code 122;
endop
rules for Pid
declare x:Pid;
bodies // should be an alias, but it does not work
  []    stdin   => 0    end
  []    stdout  => 1    end
  []    stderr  => 2    end
end of rules
rules for bool
declare x:Pid; e:int;
bodies
  []    iserr(Error(e)) => true         end
  []    iserr(x) => false               end
end of rules
rules for int
declare x:Pid; e:int;
bodies
  []    errno(Error(e)) => e            end
end of rules

strategy IO 
   builtin(1)
end of strategy

end of module
\end{verbatim}
}

\subsubsubsection{io\_aux}

The module $io\_aux$ defines the operation $read$ and $write$.
It should be imported with two arguments, where $X$ is the name of
ELAN sort and $IX$ is the module which defines the sort $X$.
In many cases (at least in the ELAN library), a sort $X$ is defined
in the module $X$. That is why we have defined module $io.eln$
(belove), which simplifies the importation.
%.. io_aux
{\scriptsize
\begin{verbatim}
module io_aux[X,IX]
// !!! not implemented in the compiled version !!!
// X should be a sort, IX module with sort X
import global IX bool int stdio string;
op global
  write(@,@)     : (Pid X) X   pri 200 code 119;     // write to a file/output to a pipe
  read(@)        : (Pid) X     pri 200 code 120;     // read from a file/get from a pipe

  iserror(@)     : (X) bool;
  errno(@)       : (X) int;
 local
  Error(@)       : (int) X     pri 200 code 123;
endop
rules for bool
declare x:X; e:int;
bodies
  []  iserror(Error(e)) => true         end
  []  iserror(x) => false               end
end of rules
rules for int
declare x:X; e:int;
bodies
  []  errno(Error(e)) => e              end
end of rules
\end{verbatim}
}
%
$write(pid,x)$ writes a term $x:X$ to a file $pid$ or to an output pipe of the
process $pid$ and in the successful
case, it returns the same term $x$. If an error is occured, the result of $write$ is
a term $Error(e)$, where $e$ is a index of error.
$read(pid)$ reads a term of the type $X$, in the error case, it returns a term $Error(e)$.

\bigskip 

{\scriptsize
\begin{verbatim}
module io[X]
import global io_aux[X,X];
end of module
\end{verbatim}
}
}
}
%----------------------------------------------------------------
\subsection{Selectors and modifiers}

The following proposition has been extracted from a `pub'-mail.

{\scriptsize
Grace a l'inspiration de Pierre E. j'ai fait une petite extension d'Elan.
Cette extension generalise la syntax du profile de vos operateurs, qui
normalement, contient les sortes du domaine et co-domaine.
Maintenant, vous {\bf pouvez} pre-fixer les sortes avec des identificateurs
(les noms de champs):
\begin{verbatim}
operators global
  [@,@] : (first:int second:bool) Pair;
end        ^^^^^     ^^^^^^
\end{verbatim}
et vous allez obtenir (pour chaque champ nome') un selecteur et un modificateur
cad:
\begin{verbatim}
  @ .first : (Pair) int;               @ .second  : (Pair) bool;
  @[.first<-@] : (Pair int) Pair;      @[.second<-@] : (Pair bool) Pair;
\end{verbatim}
avec la semantique:
\begin{verbatim}
 (x,y).first => x              
 (x,y).second => y
 (x,y)[.first<-z] => (z,y)     
 (x,y)[.second<-w] => (x,w)
\end{verbatim}

La syntaxe du selecteur et modificateur est fixe' et la semantique
est definie par deux regles de reecriture (assez evidentes).
Donc, pas de choses magiques, eventuelement, on va remplacer les
regles par des built-ins, si on aime cette extension (justement pour
la raison d'efficacite').
Par contre, je pense, que l'utilisation de cette extension va changer
le style qu'on utilise pour vos 'medium-size' elan applications.

Bon, pour illustration, ma petite 'self-explaining' application:
module sel

\begin{verbatim}
import global int bool; end
sort Pair; end
operators global
  [@,@] : (first:int second:bool) Pair;
  f(@)  : (Pair) Pair;
end

rules for Pair
p:Pair; r:int;
bodies
  []    f(p) => p[.first<-p.first+5][.second<-true] if p.first < 5      end
end
\end{verbatim}
}
%----------------------------------------------------------------
\subsection{Switch construction}

{\scriptsize
The syntax of the {\bf switch} construction is the following:
\[\begin{array}{ll}
 l \Rightarrow & {\bf switch~} \\
            & \quad whifs \\
            & \quad {\bf ~case~} c_1 {\bf ~then~} r_1 \quad whifs_1 \\
            & \quad \dots \quad \dots \quad \dots \quad \dots\\
            & \quad {\bf ~case~} c_n {\bf ~then~} r_n \quad whifs_n \\
            & \quad {\bf ~otherwise~}  \quad ~r_0 \quad whifs_0 \\
            & {\bf ~end~} \\
 {\bf ~end~}
\end{array}\]
where the sequence of wheres and ifs is defined as follows:
\[
 whifs     ::= \{ {\bf ~where~} affectation  ~|~ {\bf ~if~} c \}^{*}
\]

\noindent
The semantics of the {\bf switch} construction is the following:
\[\begin{array}{lll}
  l \Rightarrow & r_1 \quad whifs \quad {\bf ~if~} c_1 & whifs_1 \\
  l \Rightarrow & r_2 \quad whifs \quad {\bf ~if~} c_2 {\bf ~if~} not(c_1) & whifs_2 \\
            & \quad \dots \quad \dots \quad \dots \quad \dots\\
  l \Rightarrow & r_n \quad whifs \quad {\bf ~if~} c_n {\bf ~if~} not(c_{n-1}) {\bf ~if~} not(c_1) 
                         & whifs_n\\
  l \Rightarrow & r_0 \quad whifs &  whifs_0\\
\end{array}\]

\noindent
The construction {\bf ~switch~} is recursive.

See Castro's csp-solver for even nested switch constructions.
}

%----------------------------------------------------------------
\subsection{Construction where with a pattern}

{\scriptsize
The syntax of the {\bf where} statement with patterns is the following:
\[
  {\bf ~where~} (sort)pat ~:= ~[~ ()term ~|~ (s)term ~|~ [s]term ~]
\]
where the pattern $pat$ if of the type $sort$.

\noindent
The pattern $pat$ of the sort $s$ is defined as follows:
\begin{itemize}
  \item $pat$ is a variable of the type $s$, either
  \item $pat = f(p_1, \dots, p_n)$, where the symbol $f$ has the 
    profile $f:(s_1 \dots s_n)~s$ and the pattern $p_i$ has the type $s_i$,
    for any $i = 1..n$.
\end{itemize}
Even non-linear, and/or AC-patterns are treated.

There exists an example instead of some proposition:

\begin{verbatim}
module quick
import global int list[int]; end
sort pair; end
operators global
  sort(@)               : (list[int]) list[int];
  pivot(@,@,@,@)        : (int list[int] list[int] list[int]) pair;
  [@,@]                 : (list[int] list[int]) pair;
end

rules for list[int]
x       : int;
xs,s,l  : list[int];
global
  []  sort(nil) => nil                                  
  end
  []  sort(x.xs) => append(sort(s),x.sort(l))           
                where (pair)[s,l]:=()pivot(x,xs,nil,nil)
  end
end
rules for pair
p, x    : int;
xs,s,l  : list[int];
global
  []  pivot(p,nil,s,l) => [s,l]                 
  end
  []  pivot(p,x.xs,s,l) => switch
                             case p < x then pivot(p,xs,s,x.l)
                             otherwise pivot(p,xs,x.s,l)
                           end
  end
end
end
\end{verbatim}}

%----------------------------------------------------------------
\subsection{meta\_apply and its family}

\marginpar{See the library modules Meta\_apply and Meta\_strat}

    The built-in symbol \verb|meta_apply(@,@)| with the profile
    $(Strategy~X)~ X$ applies
    an ELAN strategy on a term of the sort $X$.
    It returns a result of the same sort.
    Roughly saying, the ELAN construction
    \verb@where y:=(META)meta_apply(S,t)@
    has the same effect as
    \verb@where y:=(s)t@.
    The difference between two of them comes from the relation between
    two objects (strategies) 
    $s$ and $S$. While $s$ is an ELAN built-in strategy
    (a language construction), $S$ is a term-representing an ELAN
    built-in strategy.
    This means that $s$ is a language object parsed by the
    ELAN system and transformed into its internal form
    (which may be
    an internal data structure, or a compiled c++ code).
    However, $S$ represents an ELAN term constructed
    in the run-time, and it adopts the representation of terms.
    The operational semantics of
    \verb|meta_apply(S,t)| embody the interface between these two levels.
    More precisely, \verb|meta_apply(S,t)| transforms the run-time
    representation of the strategy $S$ into its internal form $s$,
    then, it applies the obtained strategy $s$ on the term $t$. The results of
    this applications are results of the \verb@meta_apply(S,t)@.

    There exists also a sightly modified construction
    \verb|meta_apply(@,@,@) : (Strategy X int) X|
    meaning that 
    \verb|where y:=(META)meta_apply(s,t,n)| gives $n$-th solution
    of the application of $(s)t$ (if it exists).

    Values of the sort $Strategy$ are constructed by several constructors
    defined in the module \verb@Meta_strat@,
    which mimics the syntax of ELAN built-in strategies.
    We show some of them:
    {\small
    \begin{verbatim}
    sort Labels Strategies Strateg Strategy;
    operators global
     @        : (Strateg) Strategy;
     @ @      : (Strateg Strategy) Strategy; // list of sub-strategies

     id       : Strateg;                     // several primitive
     fail     : Strateg;                     // strategies

     dk(@)    : (Labels) Strateg;            // non-deterministic strategies
     dk(@)    : (Strategies) Strateg;           
     dc(@)    : (Labels) Strateg;            // non-deterministic strategies
     dc(@)    : (Strategies) Strateg;         

     repeat*( @ ) : (Strategy) Strateg;      // iterators

     @        : (Label) Strateg;             // call a sub-strategy
     \end{verbatim} }
     The signature describing term of the sort $Strategy$ is chosen
     as close as possible to the syntax of the ELAN built-in strategies.
     There are two main reasons for it:
         to do not introduce another strategy formalism, and,
         to make the transformation of $S:Strategy$ into $s$ as simple as possible
     However, there is no principal
     obstacle to enrich (or, modified) the structure of the sort $Strategy$
     except the fact, that the transformation of a sort $Strategy$ into its
     internal form should be also adapted to this extension.

     In the compiled version of ELAN,
     the transformation $S$ into its internal form $s$ means 
     a compilation of $S$.
     Effectively, to be able to support the full semantics of 
     \verb@meta_apply@ in the compiled version, we would have to
     recompile a strategies $S$ in the run-time. This explicitely
     implies the presence of the ELAN compiler in the run-time code
     of an ELAN program, which is not (up to now) our intended solution.
     That is why, there is a restriction of
     \verb@meta_apply@
     for the compiled version of ELAN, 
     which means, 
     that the strategy $S$ should 
     be a name of already defined (i.e. in compile time) strategy.
    
     The symbol \verb|meta_apply(@,@)| is an ELAN symbol (or, a constructor),
     which is evaluated due to an ELAN built-in strategy $META$.
     In the system ELAN, 
     each reduction (normalisation) 
     of a term into its normal form is deterministic,
     and the only source of the non-determinism comes from
     strategies.
     Thus, following this concept of ELAN,
     the built-in symbol
     \verb|meta_apply(@,@)| cannot be defined as a built-in functional
     symbol
     (and, applied as \verb@where y:=()meta_apply(S,t)@),
     because it would return only one solution (its normal form).
     As we have explained above, according the strategy $S$
     (if it is non-deterministic), there may exist more than one solution
     of \verb|(s)t|. Thus, the application of \verb|meta_apply| has
     a form of a reduction the strategy $META$ over the term
     \verb|meta_apply(S,t)|.

An example of using \verb|meta_apply| is a definition of the symbol
\verb|set_of(s,t,n,m)| returning
a set of $m$ solutions of the application $(s)t$ from $n-$th solution:
%
{\small
\begin{verbatim}
  set_of(@,@,@,@) : (Strategy X int int) list[X]; 

  []    set_of(S,t,n,0) => nil                    
  []    set_of(S,t,n,m) => res.set_of(S,t,n+1,m-1)
               where res:=(META)meta_apply(S,t,n) 
  []    set_of(S,t,n,m) => nil                    
\end{verbatim}}
%
The first rule terminates the recursion when all $m$ solutions
have been already found.
The second rule asks and collects 
next solutions of $(S)t$. If they do not exist, the third
rule terminates the recursion either.

Semantically, the definition of the symbol \verb|set_of|
by this system of three rules is clear.
From the point of view the complexity,
the implementation of \verb|set_of| is not ideal, because it is
quadratic on the size of $m$. The problem comes from the fact,
that to be able to find out $m$-th solution of the application
$(s)t$, we have to compute-and-skip the first $m-1$ solutions.
Thus, this solution can be viewed as a nice specification of
an eventual c++ implementation, which is not available yet.
The principal problem of the
re-implementation of \verb|set_of|
in C++ is that we need to build-in the sorts list[X].

A nice example of using the symbol \verb|set_of| is 
the definition of \verb|setof| for defined strategies.
The signature is $(\Strat{X \mapsto Y}~X)~list[Y]$ and 
its definition:
{\small
\begin{verbatim}
  []    setof(s,t,n,m) =>  set_of(eval,[s]t,n,m) 
\end{verbatim}}
The key idea of this definition is that whenever
we need to evaluate an application
of a defined
strategy $s:\Strat{X \mapsto Y}$ on a term $t$,
we construct an application term $[s]t$,
over which the principal ELAN strategy $eval$ of the
meta-interpreter of defined strategies is executed.

%----------------------------------------------------------------
\subsection{Modifications of built-in strategies}

\begin{itemize}
  \item strategies {\bf id}, and {\bf fail},
  \item strategy {\bf repeat+},
  \item strategy {\bf normalise},
\end{itemize}

There are small examples showing the force of the new 
strategy primitive  $normalise$:

\begin{itemize}
\item {\bf lambda terms normalised by $\beta\eta$ rules}
{\scriptsize\begin{verbatim}
module m
import  global  int;
        local   bool eq[lterm];
end
sort lterm;
end
operators global
        @               : (int) lterm;
        (@ @)           : (lterm lterm) lterm;
        la @            : (lterm) lterm;
end
rules for lterm
 v, w   : int;
 M, N, Q        : lterm;
local
  [beta](la M N) => repl(1,M,N)                         end
  [eta] (la M 1) => M if not free(1,M)                  end
end
strategies for lterm
global 
  [ss]  normalise(dc(beta||eta))
  end
end
\end{verbatim}}
\item {\bf derivation of polynoms}
{\scriptsize\begin{verbatim}
module poly4[Vars]

import global int Vars eq[variable] identifier list[identifier];
end
sort variable poly;
end
operators global
  FOR EACH Id:identifier SUCH THAT Id:=(listExtract) elem(Vars) :
 { Id : variable; } 

  @             : ( variable ) poly;
  @             : ( int ) poly;
  @ + @         : ( poly poly ) poly  assocRight pri 1 (AC);
  (@ + @)       : ( poly poly ) poly  alias @ + @:;
  @ * @         : ( poly  poly) poly  assocRight pri 2 (AC);
  (@ * @)       : ( poly  poly) poly  alias @ * @:;
  deriv(@,@)    : ( poly variable) poly;

  pp(@)         : (poly) poly;
end

rules for poly
  P                     : poly;
  p1, p2, p3, p4        : poly;
  x, y                  : variable;
  n, n1, n2, n3         : int;
local

  []  P+0               => P                            end
  []  P*0               => 0                            end
  []  P*1               => P                            end
  []  n1*n2             => n where n:=() n1*n2          end

  [factorize]  P + p1+p1         => P + 2*p1            end

  [factorize]  P + n1*p1 + n2*p1 => P + n3*p1
                where n3:=()n1+n2                       end

  [factorize]  P + (p1*p2) + (p1*p3) => P + p1*(p2+p3)  end

  [expand]  P + (p1+p2)*(p3+p4) => P + p2*p4 + p2*p3
                                     + p1*p4 + p1*p3    end

  [expand]  P + (p1+p2)*p3 => P + p2*p3 + p1*p3         end

  []  deriv(x,x)        => 1                            end
  []  deriv(y,x)        => 0    if $x != $y             end
  []  deriv(n,x)        => 0                            end
  []  deriv(p1+p2,x)    => p3
                where p3:=(simplify) deriv(p1,x)+deriv(p2,x)
  end
  []  deriv(p1*p2,x)    => p1*deriv(p2,x) + deriv(p1,x)*p2
  end
end

strategies for poly
global
  [simplify]
    normalise( dc(expand) ) ; normalise( dc(factorize) )
  end
end
\end{verbatim}}
\end{itemize}
%----------------------------------------------------------------
\subsection{Dynamic typing}

The following proposition has no been updated for the version elan.test4.

{\scriptsize
In this section, we explore an idea based on a simple sort
construction called {\em dynamics} \cite{DYNA-ML}, or
{\em dynamic typing} \cite{DYNATYPING,DYNAMIC}, which allows
elegantly work with sorted terms in a non-sorted manner. 
It allows also to  define functions that essentially require
run-time type-checking in statically typed language.

A weakness of static typing can be eliminated by the
construction of {\em objects with dynamic types}
(or dynamics, for short) \cite{MODULA3}, or by the introduction 
of an {\em advanced polymorphic system} \cite{MilnerTofteHarper91}.
Both methods were combined in \cite{3LOGY}, where the concept
of dynamics is generalised into {\em dependent data types}.
Another approach combining dynamics and the parametric polymorphism
is studied in \cite{DYNA-ML} from the point of the {\em type-inference}.

The definition $type~Any = t:Type,~ x:t$ introduces a
sort $Any$ (known also as $dyn$, \cite{DYNA-ML,DYNATYPING}), 
which is the simplest example of a dependent data type.
The first part of the pair $Any$ indicates the sort of the
second part. 
This means that the type indicator $t$ ranges over a domain
of all types, and the type of $x$ depends on value of $t$.
There are different approaches with respect to the construction
of the set of all types.
The situation becomes simpler moving the idea of dynamics 
into a framework of the first-order logic, 
where a set of all sorts $\SS$ is finite.

The sort $Any$ could be defined in the parametric module $any[X]$:
%
{\small
\begin{verbatim}
module any[X]
sort Any;
op global
  (X)@   : (X) Any;
endop
end of module
\end{verbatim}}\noindent
%
where the argument $X$ of this module is supposed to be a name of a sort
(like $int$, or $list[int]$). This module defines also one
{\em embedding} coercion wrapping up terms of the sort $X$ into terms of
the sort $Any$, in a such way, that if $t:X$, then $(X)t:Any$.
A projection $(X^{-1})@  ~:~ (Any) \mapsto X$ 
from the sort $Any$ into $X$ can be also defined  with the following 
semantics:
$(X^{-1})(X)t \rewrite t$.
 
In this section, we show that the construction of dynamics gives 
a possibility:
\begin{itemize}
 \item
   to combine advantages of statically-typed non-polymorphic 
   many-sorted languages, like ELAN,
   with flexibilities of untyped (or mono-typed) systems, like LISP,
 \item
   to define functions in a 'ad-hoc' polymorphic style,
 \item
   to introduce a set of term-handling constructors and and de-constructors,
 \item
   to write polymorphic term-traversal strategies.
\end{itemize}

\subsubsection{'ad-hoc' polymorphic functions}

A simple example of 'ad-hoc' polymorphic function is a function $export$
converting (or, exporting) terms into a stream of lexems, i.e. list of lexems.
The 'ad-hoc' polymorphism of $export$ means that it works
for terms of different types, and terms of different type are treated
in different ways.

Supposing, that there are only two types of lexems: 
integers and strings,
they can be represented as $Any$-terms, e.g.
$(int)n$ and $(string)s$, where $n:int$ and $s:string$. 
The profile of the function $export$ could be $(Any) \mapsto list[Any]$.
There are two elementary cases:
\[\begin{array}{ll}
  export((int)n)     & \rewrite (int)n ~.~ nil \\
  export((string)n)  & \rewrite (string)n ~.~ nil \\
\end{array}\]
and other sorts are treated by another set of rules.
For instance, the following rules for the sort $list[X]$ 
can be used:
\[\begin{array}{ll}
  export((list[X])nil)     & \rewrite (string)"nil" ~.~nil \\
  export((list[X])a.as)    & \rewrite 
       export((X)a) + (string)"." ~.~ export((list[X])as)
\end{array}\]
where the symbol $+$ stands for the $append$ over $list[Any]$.
For terms of an arbitrary type $s$, there are rules for
each constructor $f : (s_1~ \dots~ s_n) \mapsto s$ 
of the type $s$, where $t_i$ are variables of the sorts $s_i$:
\[\begin{array}{ll}
  export((s)f(t_1,\dots,t_n)) \rewrite & (string)"f(" ~.~ \\
  & export((s_1)t_1) + (string)"," ~.~ \\
  & \dots  \quad \dots \quad  \dots \quad \dots \\
  & export((s_n)t_n) ~.~ (string)")" ~.~ nil
\end{array}\]
These rules can be generated by the pre-processor of ELAN.

A definition of the inverse (or, parsing) function 
$import : (list[Any]) \mapsto Any$ is from the algorithmic
point-of-view more complicated due to a parsing algorithm.

\subsubsection{explode, implode, functor}

As we could see on the last example, the sort $Any$ allows to
define 'ad-hoc' polymorphic functions handling with terms of
an arbitrary sort. There is a small obstacle that for 
each such function and each term
constructor we should define (in general) one rule defining the
behaviour of the function over that term. The solution mentioned
above says that these rules can be generated by pre-processor of ELAN.
The following solution seems to be conceptually simpler.

Basic operations with $Any$ terms are a composition and a decomposition,
which are defined by two function symbols $explode$ and $implode$. 
The symbol $explode$
destructs an $Any$-term into a list of its arguments (or sub-terms), 
which are also wrapped into the $Any$-envelope. 
To be able to construct back the initial term, we introduce
a symbol $implode$, which from a list of arguments 
(or, sub-terms $t_i$) in the $Any$-form,
and from a function symbol name $f$ constructs the term $f(t_1, \ldots, t_n)$
(also wrapped into the $Any$-form).
If $f$ is the top most function symbol of a term $t$ (i.e. $f = functor(t)$), 
then $t = implode(f,explode(t))$ holds.
As a consequence of this, we have to deal with function symbol names as terms (or values).
Therefore, we introduce a new sort $Functor$, which collects all names of function symbols in such way that if
$ f(@,\ldots,@) : (s_1 \ldots s_n) \mapsto s \in \Sigma$, 
then the following symbol is introduced:
$ f(s_1 \ldots s_n)  :  Functor$.

The following module $Any$ introduces signatures of the symbols $explode$ and
$implode$, the sort $Functor$, and a function symbol $functor$ which extracts
a function symbol name from an $Any$ term.
%
{\small
\begin{verbatim}
module Any
import global list[Any];
sort Any Functor;
op global
  functor(@)    : (Any) Functor;
  explode(@)    : (Any) list[Any];
  implode(@,@)  : (Functor list[Any]) Any;
endop
end of module
\end{verbatim}}\noindent
%
Let us illustrate these ideas on the following example: \\
%
By importing $any[int]$ and $any[tree[int]]$, where $tree[X]$ is defined 
as follows:
%
{\small
\begin{verbatim}
op global 
  empty         : tree[X];
  leaf(@)       : (X) tree[X];
  node(@,@,@)   : (tree[X] X tree[X]) tree[X];
endop
\end{verbatim}}\noindent
%
we can destruct an $Any$ term 
\[
(tree[int])node(leaf(3),4,node(leaf(5),6,empty))  : Any
\]
by $explode$ into a list of $Any$ terms, i.e.
\[\begin{tabular}{l}
        explode((tree[int])node(leaf(3),4,node(leaf(5),6,empty)) \\
        \qquad \rewrite
        (tree[int])leaf(3)  . 
        (int)4  . 
        (tree[int])node(leaf(5),6,empty)  . 
        nil
\end{tabular}\]
To get back, we use the symbol $implode$ applied on this list with a function
symbol name $node(tree[int],int,tree[int])  :  Functor$, and we obtain:
\[\begin{tabular}{l}
        implode(
        node(tree[int],int,tree[int]), \\
        \qquad (tree[int])leaf(3)  .
        (int)4  .
        (tree[int])node(leaf(5),6,empty)  .  nil) \\
        \qquad \rewrite 
        (tree[int])node(leaf(3),4,node(leaf(5),6,empty))
\end{tabular}\]
%
% HELENE: devepopper 
This feature allows us to include typed terms using type flags 
(or type decorations) into mono-sorted $Any$-terms. We can also
make projections of $Any$-terms into a particular sort $X$, which
is a run-time test of a type flag (it can fail).
We can define various symbols and strategies over this sort $Any$,
which can be viewed (from the point of view of many-sorted systems)
as polymorphic symbols or strategies. \\
% HELENE: devepopper  end
%
From the implementation point of view, the semantics of 
$explode$, $implode$ and $functor$ are defined by three sets
of rewrite rules. 
The only trick is that these rewrite rules 
are generated and imported automatically,
whenever the user declares 
(by the importation of the module $any$)
the usage of this feature.
This is the only point making from this construction a built-in 
feature.

From the users point of view, the user should 
only declare importing of modules $any[X]$ (for various sorts $X$) and
the system automatically generates (and imports) semantical rules
defining the symbols $explode$, $implode$ and $functor$ and
it also defines appropriate function symbol names. 
The importations of the module $any$ can be viewed in terms of the
previous section as strategy declarations.
% HELENE: what are strategy declarations

One triple of rewrite rules produced for the symbol $node(@,@,@)$ is
an example of symbols and rules automatically generated:
%
\[\begin{tabular}{l}
  explode((Any)node(x,y,z)) \rewrite 
        (tree[int])x  .  (int)y  .  (tree[int])z  .  nil \\
  implode(node(tree[int],int,tree[int]),~
        (tree[int])x  .  (int)y  .  (tree[int])z  .  nil) \\
        \qquad \rewrite
        ((Any)node(x,y,z)) \\
  functor((Any)node(x,y,z)) \rewrite 
        node(tree[int],int,tree[int])
\end{tabular}\]
%
where $x, z : tree[int]$ and $y : int$ are variables. These rules
can be correctly parsed if the following function symbol name is declared:
\[
node(tree[int],int,tree[int]) : Functor
\]
Similar rewrite rules and symbol definitions are also automatically produced
for other constructor of the sort $tree[int]$, and in general, 
they are produced for any 
function symbol $f : (s_1 \ldots s_n) \mapsto s \in \Sigma$, 
if the user has already 
imported modules $any[s_i]$ and $any[s]$, for any $i=1..n$.

\bigskip

Using these general $Any$-term constructors and destructor 
we can easily define
function accessing $n$-th argument of an arbitrary term $a$ ($n-th~arg(a)$),
or function replacing $n$-th argument of $a$ by $b$ ($a[n-th b]$):
\[\begin{array}{ll}
    @-th~arg(@)      & : (int~Any) \mapsto Any   \\
    @[@-th~@]   & : (Any~int~Any) \mapsto Any
\end{array}\]
where for the definition of their semantics we use library list-indexing 
functions:
\[\begin{array}{ll}
    n-th~arg(a)      & \rewrite n-th~elem(explode(a)) \\
    a[n-th~b]        & \rewrite implode(functor(a),
                          replace~ n-th~elem(explode(a))~by~b)
\end{array}\]

\noindent
Similarly, we can re-define ELAN's built-in library functions
$occur~a~in~b$, or $replace~a~in~b~by~c$.

\subsubsection{Polymorphic strategies}

Up to now (in this section), we have illustrated the idea, which allows to
write `polymorphic' rewrite systems. Now, we concentrate ourself to the definition
of `polymorphic' strategies, i.e. strategies working over an arbitrary user defined sort. 
There is a strong motivation for this kind of extension, otherwise
we are not able to define strategies (like `left-most inner-most')
which work over any sort. What we are able to do, 
is to define those general strategies
for any particular sort, for example $tree[int]$. This means that
in a such definition, it is explicitly expressed that it runs
over a particular sort,  because it contains their particular term 
constructors and destructor.
However, our intention is to define a polymorphic version of 
the `left-most inner-most' strategy, which runs over any sort, i.e. over
$tree[int]$ and over lambda terms as well.

The following module defines a strategy $lis(s) : (\Strat{X}) \mapsto \Strat{Any}$
(an abbreviation for `left-most inner-most'), which applies
an argument strategy $s : \Strat{X}$ on the
left-most inner-most redex of the sort $X$.
%
{\small
\begin{verbatim}
module lis[X]
import global Any list[Any] strat[X] strat[Any] strat[list[Any]];
op global
  lis(@)  : (<X>) <Any>;
 local
  map(@)  : (<X>) <list[Any]>;
  mis(@)  : (<X>) <Any>;
endop
strategy for <Any>
declare s   : <X>;     a   : Any;    as  : list[Any];
body
  []  lis(s)    ==> case(mis(s),(X)s)                   end
  []  [mis(s)]a =>  implode(functor(a),as) 
                        where as:=[map(s)]explode(a)    end
end of strategy
strategy for <list[Any]>
declare s  : <X>;
body
  []  map(s) ==> case(cons(lis(s),id), cons(id,map(s))) end 
end of strategy
end of module
\end{verbatim}}\noindent
%
This definition introduces two auxiliary strategies,
$map$ and $mis$. 
The definition of the strategy $mis$ is temporary 
(cf. section \ref{let}),
and the definition of $lis(s)$ will be later simplified
by introducing a new strategy constructor $let-in$.
The strategy $mis(s)$ assigns to the variable $as$
a  list of sub-terms of $a$ over which the $map(s)$ strategy is
applied. Eventually, we obtain a result of this application and
the final result (an $Any$ term) is constructed by $implode(functor(a),as)$.
Otherwise, the second sub-strategy in the $case$ expression, i.e.
the strategy $(X)s  :  \Strat{Any}$, is applied.
This strategy expression is little bit tricky, however for its explanation 
it is sufficient
to recall that $(X)@  :  (X) \mapsto Any$ is a function symbol defined in the
module $any[X]$. That is why, there is also an elementary strategy symbol
$(X)@:(\Strat{X}) \mapsto \Strat{Any}$, 
which allows us to combine the strategy 
$s:\Strat{X}$
into 
the strategy
$(X)s  :  \Strat{Any}$. The operational meaning of this strategy expression is:
when it is applied on an $Any$ term, it destructs the $Any$ term into a term of the 
sort $X$ (which is, in fact, a test of the type flag)
and then, it applies the strategy $s$ on the obtained
term. Finally, it constructs back a result term
from a result of the strategy application 
(which has a sort $X$) by its wrapping by $(X)@$ to the form of an $Any$ term. 
The resume is that this notation might be seen little bit
difficile for the first reading, however, there is noting hacked behind it. 

The second auxiliary strategy $map(s)  :  \Strat{list[Any]}$ tries to
apply the strategy $lis(s)$ to the left-most element of the input list.
Actually, it applies $lis(s)$ on the head, and if it fails,
$map(s)$ continues on the tail of this list. If the input list does not contain
any element, on which the strategy $lis(s)$ is applicable, $map(s)$
fails also. This causes the invocation of the strategy $(X)s$ (as it was described
in the previous paragraph).

If we swap two strategy expressions in the $case$ of $lis(s)$
%
{\small
\begin{verbatim}
  []  lis(s) ==> case(s,mis(s))                         end
\end{verbatim}}\noindent
%
we obtain an analogical `left-most outer-most' polymorphic strategy
(which, for better understanding, should be renamed to $los(s)$).

Now, we have the first definitions of $lis$ and $los$ strategies.
% If we import the module $lis[int]$ and having any strategy $s$
% of the type $\Strat{int}$, we can apply $lis(s)$ on the $tree[int]$.
In this paragraph, we would like to stress the fact, that
the argument $s$ of $lis(s)$ in its definition should by
``typed-strategy'', i.e. a strategy of the type $\Strat{X}$ for certain
sort $X$. We also want to show another definitions of the strategies
$lis$ and $los$, which work with argument strategies $s$ of the sort 
$\Strat{Any}$.
As an example of a such strategy, let us look at the following 
{\em artificial} strategy $swap : \Strat{Any}$:
%
{\small
\begin{verbatim}
strategies for <Any>
declare x,z     : tree[int];         y       : int;
bodies
  []    [swap] (int)y => (int)y+1                                end
  []    [swap] (tree[int])node(x,y,z) => (tree[int])node(z,y,x)  end
end of strategies
\end{verbatim}} \noindent
%
This strategy is applicable to any $Any$ term, and in the case of
$int$-term, i.e. $(int)n$, where $n:int$, 
it increments its value ($n$) by $1$.
In the case of a $tree[int]$ term, 
it swaps left and right subtrees of $node$, if there are any subtrees.
Otherwise, it fails.
This could be
a meaningless example of an 'ad-hoc' polymorphic strategy, which can be
passed as an argument to the following definition of $lis$. 
The redefinition of $lis$ and $los$ is in the following 
non-parametric module $lis1$:
%
{\small
\begin{verbatim}
module lis1
import global Any list[Any] strat[Any] strat[list[Any]];
op global
       lis(@)   : (<Any>) <Any>;
       los(@)   : (<Any>) <Any>;
    local
        map(@)  : (<Any>) <list[Any]>;
        mis(@)  : (<Any>) <Any>;
endop
strategy for <Any>
declare s   : <Any>;
        a   : Any;         as  : list[Any];
body
  []  lis(s) ==> case(mis(lis(s)),s)                    end
  []  los(s) ==> case(s,mis(los(s)))                    end
  []  [mis(s)]a => implode(functor(a),as) 
                           where as:=[map(s)]explode(a) end
end of strategy
strategy for <list[Any]>
declare s : <Any>;
body
  []  map(s) ==> case(cons(s,id), cons(id,map(s)))      end
end of strategy
end of module
\end{verbatim}}\noindent
%
Now, we can call a new strategy $lis(s)$ parameterised by a strategy 
$s:\Strat{Any}$.
If the strategy $s$ has a type $\Strat{Y}$, we can call
the new version of $lis$ with a strategy $(Y)s$, i.e. $lis((Y)s)$,
because, as we have shown, $(Y)s : \Strat{Any}$.
This small example shows a way how to convert a $\Strat{Y}$ strategy into 
a $\Strat{Any}$ strategy.
Later (cf. section \ref{inline}), 
we show a conversion of an $\Strat{Any}$ strategy into a $\Strat{X}$ strategy.
}

%----------------------------------------------------------------
\subsection{HO-extension}

The following proposition has no been updated for the version elan.test4.

{\scriptsize
In this section, we illustrate another simple extension of ELAN,
which opens ELAN's doors to typical second-order applications
known from the functional programming world.
We propose a way how one can elegantly define and use the
second-order functions,
like {\em uniform traversal combinators} \cite{COMBIN-CADE-11},
or
{\em catamorphisms},
{\em anamorphisms},
{\em hylomorphisms} and
{\em paramorphisms} \footnote{
in functional programming community familiarly called
bananas, lenses, envelopes and barbed wires} \cite{Bananas}.

The stress in this section is not put on the fact that we use
the first-order rewriting logic to mimic a manipulation with
functional objects, but on the fact, that a lot of things we can
do elegantly. There is no doubt that rewriting logic is powerful
enough to simulate $\beta\eta$-normalisation of $\lambda$-terms,
however, this is not our goal.
All functions defined in this section are weakly motivated
examples of the second-order functions well-known from
the functional programming folklore. We do not advocate, that
they represent functions, which the ELAN's programmer needs
every day. \\
To simplify our presentation,
all examples (of combinators and morphisms) in this section,
we illustrate on a recursive type $list[X]$.
Extensions for general types are out of the scope of this section.

The paper \cite{COMBIN-CADE-11} introduces a notion of the
{\em uniform traversal combinator} representing a pattern of
the recursion.  Due to a crucial result,
that any function has a unique representation as a
traversal combinator, this paper shows that these combinators
can substantially help the theorem proving process by
eliminating the need for induction.

A {\em list traversal combinator} is a function:
\[
tclist : (Y,~ (X,list[X],Y) \mapsto  Y,~ list[X]) \mapsto  Y
\]
defined as follows:
\[\begin{array}{ll}
 tclist ( fnil, fcons )( nil )  &= fnil \\
 tclist ( fnil, fcons )( a.as ) &= fcons(a,as,tclist(fnil,fcons)(as) \\
\end{array}\]
Having a function $rev : (X, list[X], list[X]) \mapsto  list[X]$ defined as
$rev(a,as,l) \rewrite append(l,a.nil)$,
we can reverse a list using this list traversal combinator:
\[
    reverse(x) = tclist(nil, rev)(x)
\]
Definitions of different $list$ traversal functions, 
like $length$, $append$, etc., 
using the list traversal combinator $tclist$ is a 
typical functional exercise.
The proclaimed goal of this section is to show their analogical
definitions in ELAN --- in the first-order framework.
The problem is the second (functional) argument of $tclist$,
and
there is no strighforward way to re-define $tclist$
using the first-order framework of rewriting logic.

We introduce a new sort $\Symb{(X,list[X],Y) \mapsto Y}$ which is 
a sort of all function symbols (from $\Sigma$) of the profile
$(X, list[X], Y) \mapsto  Y$.
In our example, there exists at least one,
$rev : \Symb{(X,list[X],Y) \mapsto Y}$.
We introduce also one operation of application of a function symbol 
$p \in \Symb{(X,list[X],Y) \mapsto Y}$ to a triple of its arguments:
\[
~@~ \lban ~@,~@,~@~ \rban ~:~ 
    (\Symb{(X,list[X],Y) \mapsto Y},X,list[X],Y)\mapsto Y
\]
and this operation $@\lban @,@,@ \rban$ has the following interpretation:
\[
p \lban x, xs, y \rban \rewrite p(x,xs,y)
\]
This construction allows us to re-define the function 
$tclist$ as follows:
\[\begin{array}{ll}
  tclist(fnil,fcons)(nil)  &\rewrite fnil \\
  tclist(fnil,fcons)(a.as) &\rewrite 
  fcons \lban a,as,tclist(fnil,fcons)(as) \rban
\end{array}\]
and define it $reverse(x) = tclist(nil,rev)(x)$, where
$fnil,as : list[X]$, $a:X$ and $fcons : \Symb{(X,list[X],Y) \mapsto Y}$
are first-order variables, while
$rev: \Symb{(X,list[X],Y) \mapsto Y}$ is a constant representing
an index of the function $rev: (X,list[X],Y) \mapsto Y$.

The function $tclist$ is an example of a
{\em paramorphism}%\footnote{barbed wire}
studied in \cite{Meert}.
All mentioned examples ($length$, $append$ and $reverse$) are,
however,
examples of {\em catamorphisms}%\footnote{bananas}
, because
the function $fcons$ does not depend on its second argument.
Thus, the definition of the list-catamorphisms can be simplified 
to the form:
$catalist : (Y, \Symb{(X,Y) \mapsto Y}, list[X]) \mapsto  Y$,
programmed as:
\[\begin{array}{ll}
  catalist(fnil,fcons)(nil)  &\rewrite fnil \\
  catalist(fnil,fcons)(a.as) &\rewrite 
  fcons \lban a,catalist(fnil,fcons)(as) \rban
\end{array}\]
The catamorphism $catalist$ is known in functional programming
as $foldr$.

One example of {\em paramorphism} not definable by $catalist$
could be a function computing all tail segments of a list,
i.e. $tails : (list[X]) \mapsto list[list[X]]$, which is an
instantiation of $tclist$, where $Y = list[list[X]]$ and
\[
   tails(x) = tclist(nil.nil,cons\_cons)
\]
where $cons\_cons : (X, list[X], list[list[X]]) \mapsto list[list[X]]$
is defined as follows:
\[\begin{array}{ll}
  cons\_cons(a,as,ass) = a.as.ass
\end{array}\]

The last example of a morphism known in the functional programming crowd
as {\em anamorphism}%\footnote{lenses} 
is illustrated on
the function $unfold$:
\[
  unfold : ( Y \mapsto bool, Y \mapsto Pair[X,Y], Y) \mapsto list[X]
\]
\[\begin{array}{lll}
  unfold(p, g)(b) & = nil                & \IF p(b) \cr
  unfold(p, g)(b) & = a . unfold(p,g)(c) & \WHERE [a,c]:=g(b) 
\end{array}\]
The function $unfold$ is parameterised by two functional objects.
The first one is a 'predicate' terminating the iteration, which is a 
function symbol from a sort $\Symb{Y \mapsto bool}$.
The second argument develops one element of a constructed list and 
a continuation value $c$. Thus, we instantiate the second argument
by a function symbol from $\Symb{Y \mapsto Pair[X,Y]}$.
Then, the ELAN's definition is strighforward
\[
 unfold : ( \Symb{Y \mapsto bool}, \Symb{Y \mapsto Pair[X,Y]} , Y) \mapsto list[X]
\]
\[\begin{array}{lll}
  unfold(p, g)(b) & \rewrite Nil                & \IF p \lban b \rban \\
  unfold(p, g)(b) & \rewrite a . unfold(p,g)(c) & \WHERE [a,c]:=g \lban b \rban 
\end{array}\]
%
If, for example, we instantiate the anamorphism $unfold$ as follows:
\begin{itemize}
 \item 
   $Y = Pair[list[A],list[B]]$, and $X = Pair[A,B]$, where $A,B \in \SS$,
 \item
   $ p : ( Pair[list[A],list[B]] ) \mapsto bool$, s.t. \\
   $ p([as,bs]) = (as == nil) \OR (bs == nil) $,
 \item
   $ g : ( Pair[list[A],list[B]] ) \mapsto Pair[Pair[A,B] ,Pair[list[A],list[B]] ] $, s.t.\\
   $ g([a.as, b.bs]) = [~ [a,b],~ [as,bs]~ ] $
\end{itemize}
we obtain a function $zip(xs,ys) = unfold(p,g)([xs,ys])$
zipping two lists into one, i.e.
$zip([1.2.3.nil,a.b.nil]) = [1,a].[2,b].nil$.

\bigskip

\noindent
We have shown a small palette of different second-order
examples. Now, we generalise the constructions of a 
{\em system of symbol sorts}. \\
A symbol sort $\Symb{(s_1 \ldots s_n) \mapsto s}$ 
is a sort of all function symbol names $f(s_1, \ldots, s_n)$ of
function symbols $f : (s_1 ~ \ldots ~ s_n) \mapsto s$ from $\Sigma$.
For example, a sort $\Symb{(int) \mapsto bool}$ 
collects all symbol names of function symbols of
the profile $(int) \mapsto bool$, i.e. names of all integer predicates. 
For example, having a test of prime numbers $prime(@): (int) \mapsto bool$,
the symbol $prime(int)$ has a type $\Symb{(int) \mapsto bool}$.
We have slightly simplified examples in the first part of this section,
because symbol names were represented by function symbols $f$,
e.g. $primes$,
however, the definition above introduces symbol names as 
$f(s_1, \ldots, s_n)$, e.g. $primes(int)$. The reason for it is
to remove ambiguities.

The sort $\Symb{(s_1 \ldots s_n) \mapsto s}$ 
(typed as \verb@Symbol[s_1,...,s_n,s]@) is defined in the following
(sophisticatedly parameterised) library module $symbol$:
%
{\small
\begin{verbatim}
module symbol[n, s_1,...,s_n, s]
sort   Symbol[s_1,...,s_n,s];
op global
  @ (@ {,@}_I=2...n) : (Symbol[s_1,...,s_n,s] {s_I}_I=1...N) s;
endop
end of module
\end{verbatim}}\noindent
The first side-effect caused by importing
the module $symbol[n,s_1,\ldots,s_n,s]$ is that the system generates
function symbol names for all symbols of the sort 
$\Symb{(s_1 \ldots s_n) \mapsto s}$,
i.e. of the profile $(s_1 ~ \ldots ~ s_n) \mapsto s$. 
The first argument $n$ refers to the
arity of these symbols and it must be explicitly mentioned.
For any symbol sort $\Symb{(s_1 \ldots s_n) \mapsto s}$ defined by the 
module $symbol$, the module $symbol$ defines 
one $(n+1)$-ary application operation of the profile:
\[
~@~ \overbrace{ \lban ~@, \ldots, @~ \rban }^{n} ~:~ 
    (\Symb{(s_1~\ldots~s_n) \mapsto s}~s_1~\ldots~s_n)\mapsto s
\]
The second side-effect caused by importing
the module $symbol[n,s_1,\ldots,s_n,s]$ is that 
for any function symbol name $p : \Symb{(s_1 \ldots s_n) \mapsto s}$ 
of the symbol $f : (s_1 \ldots s_n) \mapsto s$
there is automatically generated rewrite rule:
\[
p \lban x_1, \ldots, x_n \rban \rewrite f(x_1, \ldots, x_n) \quad       
                        where~x_i:s_i
\]
which allows (dynamically) assemble terms from the name of its 
function symbol
$p:\Symb{(s_1~\ldots~s_n)\mapsto s}$ and their sub-terms $x_i:s_i$.

Finally, we illustrate these ideas on two definitions of simple strategies.

\noindent
The following module $depth[X]$ defines a depth-first search 
strategy $depth-first(p,s)$ iterating
the strategy $s$ while a certain condition $p$ holds.
%
{\small
\begin{verbatim}
module depth[X]
import global symbol[1,X,bool] strat[X];
op global
  depth-first(@,@) : (Symbol[X,bool] <X>) <X>;
endop
strategy for <X>
declare p : Symbol[X,bool];           s : <X>;
body
  []  depth-first(p,s) ==>
        if p(self) then
          s; depth-first(p,s)
        else
          id
        fi
  end
end of strategy
end of module
\end{verbatim}}\noindent
%
The first argument of $depth-first$ has a sort $\Symb{(X)\mapsto bool}$,
which is a domain of all function symbol names of boolean predicates 
for the sort $X$, i.e. all symbols with the profile $(X) \mapsto bool$.
The condition $p \lban self \rban :bool$ in the 
$if-then-else-fi$ strategy is assembled form 
the function symbol name $p$ and the term
$self:X$ using the application symbol and the rewrite rule
described above.

Another simple example of a strategy using this feature is the traditional  
$filter(p,s)$ strategy, which selects from a list only elements
satisfying a condition $p$, and afterwards, it applies a strategy $s$
to all accepted elements.
{\small
\begin{verbatim}
module filter[X]
import global X list[X] symbol[1,X,bool] strat[X] strat[list[X]];
op global
  filter(@,@) : (Symbol[X,bool] <X>) <list[X]>;
endop
strategy for <list[X]>
declare h, h1 : X;        t, t1 : list[X]; 
        s : <X>;          p : Symbol[X,bool];
body
  []  [filter(p,s)]nil  => nil                            end
  []  [filter(p,s)]h.t  => h1.t1 where t1:=[filter(p,s)]t
                                 where h1:=[s]h
                           if p(h)                        end
  []  [filter(p,s)]h.t  => t1 where t1:=[filter(p,s)]t
                           if not p(h)                    end
end of strategy
end of module
\end{verbatim}}\noindent
%
where the last two rules can be simplified into:
{\small
\begin{verbatim}
  []  [filter(p,s)]h.t  => switch
                             case p(h) then h1.t1
                               where h1:=[s]h
                             otherwise t1
                           end
                           where t1:=[filter(p,s)]t
\end{verbatim}}\noindent
To summarise this section, it becomes clear that we are able to
define several strategies corresponding to some higher-order
functions known from the functional programming. 
We have to stress, that each
symbol sort $\Symb{(s_1 \ldots s_n) \mapsto s}$ contains only a finite
set of symbol names (and we have no 'abstraction' construction), thus,
our extension with symbol names has nothing to do with the higher-order
rewriting\footnote{Moreover, it might be seen as an advantage}.
The proposed extension allows us to work with functions
(or function symbol names as their indexes) in the same way as it is
used in procedural languages, like c++.
}

%----------------------------------------------------------------
\subsection{Built-in data types (i.e. strings and reals)}

\marginpar{See the library module string}
There is nothing written neither about strings, nor about reals.
There is a library specification of built-ins over strings:
{\scriptsize \begin{verbatim}
   strlen(@)    : (string) int             // the length of a string
   strcat(@,@)  : (string string) string   // concatenation of strings
   @+@          : (string string) string   // alias to the concatenation of strings
   @[@]         : (string int) int         // s[i] :: selects i-th character of the string s
   @[@<-@]      : (string int int) string  // s[i<-c] :: modifies i-th character of the string s by the character c
   substr(@,@,@): (string int int) string  // substr(s,i,n) :: extracts a substring from the i-th position with n charters
   strspn(@,@)  : (string string) int      // computes the length of the maximum initial segment, see the unix function strspn
   strcmp(@,@)  : (string string) int      // compares two strings, see unix function strcmp
   string(@)    : (int) string             // conversion from a integer (character) into a string
\end{verbatim}}

and, just a small non-trivial example:

{\scriptsize \begin{verbatim}
   module moi_meme
import global int bool string; end
operators global
  moi(@)        : (int) string;
  loop(@,@,@)   : (int int int) string;
  trans(@,@)    : (int string) string;
  moi_meme      : string;
end
rules for string
i, j, k: int; 
s      : string;
local
  [] moi(00) => "module moi_meme" end
  [] moi(01) => "import global int bool string; end" end
  [] moi(02) => "operators global" end
  [] moi(03) => "  moi(@)        : (int) string;" end
  [] moi(04) => "  loop(@,@,@)   : (int int int) string;" end
  [] moi(05) => "  trans(@,@)    : (int string) string;" end
  [] moi(06) => "  moi_meme      : string;" end
  [] moi(07) => "end" end
  [] moi(08) => "rules for string" end
  [] moi(09) => "i, j, k: int; " end
  [] moi(10) => "s      : string;" end
  [] moi(11) => "local" end
  [] moi(12) => "  [] loop(i,j,k) => switch" end
  [] moi(13) => "                      case i <= j and k == 0 then moi(i)+string(10)+loop(i+1,j,k)" end
  [] moi(14) => "                      case i <= j and k == 1 then trans(i,moi(i))+loop(i+1,j,k)" end
  [] moi(15) => "                      otherwise string(0)" end
  [] moi(16) => "                    end" end
  [] moi(17) => "  end" end
  [] moi(18) => "  [] trans(i,s) => string(32)+string(32)+string(91)+string(93)+string(32)+string(109)+string(111)+string(105)+" end
  [] moi(19) => "                   string(40)+string(48+i/10)+string(48+i%10)+string(41)+string(32)+string(61)+string(62)+string(32)+" end
  [] moi(20) => "                   string(34)+s+string(34)+string(32)+string(101)+string(110)+string(100)+string(10) end" end
  [] moi(21) => "  [] moi_meme => loop(0,11,0)+loop(0,23,1)+loop(12,23,0) end" end
  [] moi(22) => "end" end
  [] moi(23) => "end" end
  [] loop(i,j,k) => switch
                      case i <= j and k == 0 then moi(i)+string(10)+loop(i+1,j,k)
                      case i <= j and k == 1 then trans(i,moi(i))+loop(i+1,j,k)
                      otherwise string(0)
                    end
  end
  [] trans(i,s) => string(32)+string(32)+string(91)+string(93)+string(32)+string(109)+string(111)+string(105)+
                   string(40)+string(48+i/10)+string(48+i%10)+string(41)+string(32)+string(61)+string(62)+string(32)+
                   string(34)+s+string(34)+string(32)+string(101)+string(110)+string(100)+string(10) end
  [] moi_meme => loop(0,11,0)+loop(0,23,1)+loop(12,23,0) end
end
end
\end{verbatim}
}

%%%%%%%%%%%%%%%%%%%%%%%%%%%%%%%%%%%%%%%%%%%%%%%%%%%%%%%%%%%%
\section{New Implementation Features}
%%%%%%%%%%%%%%%%%%%%%%%%%%%%%%%%%%%%%%%%%%%%%%%%%%%%%%%%%%%%

%----------------------------------------------------------------
\subsection{Typing of strategies}

This is about the fact, that the Marian's strategies were not typed,
while the current strategies are strictly typed. This represents
a language restriction, however, one can discuss also the possibility
of overloading strategy names over different types. The type-checking
algorithm for built-in strategies is not interesting enough for description.

Not documented yet.

%----------------------------------------------------------------
\subsection{Encapsulation of rules and strategies (i.e. local/global)}

In the latest version of \elan, the names rules and strategies could
be specified as global or local.
The only readable document is an exercise we did during one 
Tuesday session.

Not documented yet.

%----------------------------------------------------------------
\subsection{Linearisation of rewrite rules}

This feature allows to compile non-linear rules with the current
version of the compiler.

Not documented yet.

%----------------------------------------------------------------
%\subsection{Partial evaluation}

%%%%%%%%%%%%%%%%%%%%%%%%%%%%%%%%%%%%%%%%%%%%%%%%%%%%%%%%%%%%
\section{New Interface Features}
%%%%%%%%%%%%%%%%%%%%%%%%%%%%%%%%%%%%%%%%%%%%%%%%%%%%%%%%%%%%

\subsection{Compilation}

All the stuff about the compiler should be described, also
with the Pierre's option \verb@--earley@, which produces 
the independent executable file.

%----------------------------------------------------------------
\subsection{Command language}

There are two text files:
\begin{itemize}
\item a help file (from the library):
{\scriptsize\begin{verbatim}
elan -C -q ... lgi-file spc-file

  quit          - to unix
  help          - show the file $ELANLIB/help.txt
  load module   - load a 'module' 
  batch ident   - run a script file 'ident'
  run term      - 'term' may contain Q, Q1,..., R, R1,... 
                  which refer to stacks of qs and rs

  sorts type type type  - set query, result and print sorts
  startwith (str)term   - set the 'start with' term (may contain the symbol 
                          query and Q, Qi, R, Ri)
  checkwith term        - set the 'check with' term:bool (as in startwith)
  printwith term        - set the print-with term (contains result)

  qs            - queue of queries (history of the last 10)
  rs            - queue of results (history of the last 10)
                - term in qs, rs are referenciable by 'variables'
                  'Q, Q1, ..., R, R1, ...', respecting their types

  stat          - statistic
  dump          - dump all
  dump ident    - dump named rules 'ident' or strategies 'ident' or
                  rules with head symbol 'ident'
  dump int      - dump rules for a symbol with the internal code 'int'

  display int   - display 1 - show always internal forms of term
                  display 0 - tratitional form 
  trace int     - set trace level (Marian's trace value)
  break ident   - set a break-point to named rules 'ident' or to strategies 
                  'ident' or to non-named rules with head symbol ~'ident'
  break int     - set a break-point to rules with head symbol of the internal
                  code 'int'
  unbreak ident - remove break-point ...
  unbreak int   - remove break-point ...
  breaks  - list of all breakpoints
\end{verbatim}}

\item a small demostrattion file:
{\scriptsize\begin{verbatim}
elan -C -q test
        ******** ELAN version v1.17 (25/03/96.11:46) ********

           (c)   INRIA-Lorraine  &  CRIN,   1994            

 handling module int
   handling module bool
   end of bool
 end of int
 handling module test
   handling module list[int]
     handling module listCommons
     end of listCommons
   end of list[int]
 end of test
 handling module Query[int,list[int],list[int]]
 end of Query[int,list[int],list[int]]
enter command finished by ';':

run 0;
------
[] start with term :
    f(0)
[] result term:
    1.1.nil
[] result term:
    0.0.nil
[] fin
enter command finished by ';':

run 1;
------
[] start with term :
    f(1)
[] result term:
    1.0.nil
[] result term:
    1.1.nil
[] fin
enter command finished by ';':

qs;
---
Queries ... 2 : elements 
Q1      int      0
Q       int      1
enter command finished by ';':

rs;
---
Results ... 4 : elements 
R3      list[int]        1.1.nil
R2      list[int]        0.0.nil
R1      list[int]        1.0.nil
R       list[int]        1.1.nil
enter command finished by ';':

run Q+Q1;
---------
[] start with term :
    f(plus(1,0))
[] result term:
    1.0.nil
[] result term:
    1.1.nil
[] fin
enter command finished by ';'

run size(R)+size(R1);
---------------------
[] start with term :
    f(plus(size_of_int_list(1.1.nil),size_of_int_list(1.0.nil)))
[] result term:
    1.g(4).nil
[] result term:
    4.4.nil
[] fin
enter command finished by ';':

break DO1;
----------
[1]
rule DO1 [1]
 varnum = 1     match = 1       nameindex = 196
  f( VAR(0))     =>  1.g( VAR(0)).nil
end of rule
enter command finished by ';':

run 1;
------
[] start with term :
    f(1)
     trying  'DO1' on    f(1)
     'DO1' :     1.g(1).nil
     [red] :     1.0.nil
[] result term:
    1.0.nil
     still trying  'DO1' on      f(1)
     fail of 'DO1' 
[] result term:
    1.1.nil
[] fin
enter command finished by ';':

break g; unbreak DO1;
--------
[1]
rule [1]
 varnum = 0     match = 1       nameindex = -1
  g(0)   =>  1
end of rule
[2]
rule [2]
 varnum = 0     match = 1       nameindex = -1
  g(1)   =>  0
end of rule
'g' '('  @ ')'   : (int)int      pri 0 code 139;
enter command finished by ';':

run 1;
------
[] start with term :
    f(1)
 [0]  g(1)
 [1]  0
[] result term:
    1.0.nil
[] result term:
    1.1.nil
[] fin
enter command finished by ';':

breaks;
-------
[1]
rule [1]
 varnum = 0     match = 1       nameindex = -1
  g(0)   =>  1
end of rule

[2]
rule [2]
 varnum = 0     match = 1       nameindex = -1
  g(1)   =>  0
end of rule
enter command finished by ';':

checkwith small_int(query+5);
-----------------------------
enter command finished by ';':

stat;
-----
Input type : int
Output type: list[int]
Print type : list[int]
Strategy   : do
Start_with :  f( VAR(0))
Check_with :  small_int(plus( VAR(0),5))
Print_with :  VAR(0)
------------
[] statistics: time 0.033 sec
                1 nonamed rules applied,     2 tried
                3   named rules applied,     3 tried

                rule: applied   tried                 rule: applied   tried
         extractrule1       0       0          extractrule2       0       0
                  foo       1       1                   DO1       1       1
                  DO2       1       1 

enter command finished by ';':

printwith 1996.result;
----------------------
enter command finished by ';':

run 0;
------
[] start with term :
    f(0)
[] print term:
    1996.1.1.nil
[] print term:
    1996.0.0.nil
[] fin
enter command finished by ';':

startwith (fou) query.nil;
--------------------------
enter command finished by ';':

stat;
-----

Input type : int
Output type: list[int]
Print type : list[int]
Strategy   : fou
Start_with :  VAR(0).nil
Check_with :  small_int(plus( VAR(0),5))
Print_with :  1996. VAR(0)
.....

enter command finished by ';':

run 0;
------
[] start with term :
    0.nil
[] print term:
    1996.0.nil
[] fin
enter command finished by ';':

load mest[bool];
----------------
 handling module mest[bool]
   handling module list[bool]
   end of list[bool]
 end of mest[bool]

enter command finished by ';':

sorts bool bool bool;
---------------------
 handling module Query[bool,bool,bool]
 end of Query[bool,bool,bool]
enter command finished by ';':

checkwith true; printwith result; startwith (fou) true and query; stat;
----------------------------------------------------------------------

Input type : bool
Output type: bool
Print type : bool
Strategy   : fou
Start_with : (true and VAR(0))
Check_with :  true
Print_with :  VAR(0)
........
enter command finished by ';':

run false;
----------
[] start with term :
   (true and false)
[] print term:
    false
[] fin
enter command finished by ';'
quit;
.......
\end{verbatim}}
\end{itemize}


\bibliography{/users/protheo/bibtex/abbrev,/users/protheo/bibtex/protheo,/users/protheo/bibtex/ckirchne}
\bibliographystyle{plain}  
\end{document}



